\documentclass{article}
\usepackage{iclr2026_conference,times}

%%%%% NEW MATH DEFINITIONS %%%%%

\usepackage{amsmath,amsfonts,bm}

\newcommand{\newterm}[1]{{\bf #1}}
\def\figref#1{figure~\ref{#1}}
\def\Figref#1{Figure~\ref{#1}}
\def\tabref#1{table~\ref{#1}}
\def\Tabref#1{Table~\ref{#1}}
\def\secref#1{section~\ref{#1}}
\def\Secref#1{Section~\ref{#1}}
\def\eqref#1{equation~\ref{#1}}
\def\Eqref#1{Equation~\ref{#1}}

\newcommand{\R}{\mathbb{R}}
\newcommand{\E}{\mathbb{E}}
\newcommand{\Prob}{\mathbb{P}}
\newcommand{\Sphere}{\mathbb{S}}
\newcommand{\calD}{\mathcal{D}}
\newcommand{\calE}{\mathcal{E}}
\newcommand{\calN}{\mathcal{N}}
\newcommand{\calT}{\mathcal{T}}
\newcommand{\calC}{\mathcal{C}}


\usepackage{hyperref}
\usepackage{url}
\usepackage{booktabs}
\usepackage{amsmath}
\usepackage{amssymb}
\usepackage{amsthm}
\usepackage{graphicx}
\usepackage{microtype}
\usepackage{placeins}
\usepackage{flafter}
\hypersetup{hidelinks}

\setlength{\textfloatsep}{8pt plus 2pt minus 2pt}
\setlength{\floatsep}{8pt plus 2pt minus 2pt}
\setlength{\intextsep}{8pt plus 2pt minus 2pt}

\newtheorem{definition}{Definition}
\newtheorem{proposition}{Proposition}
\newtheorem{theorem}{Theorem}
\newtheorem{remark}{Remark}

\makeatletter
\renewenvironment{proof}[1][\proofname]{\par\noindent\textit{#1.}\ }{\hfill\qedsymbol\par}
\makeatother

\title{Spectral Partition Routing for Uniform Feature Dispatch in Mixture-of-Experts}

\author{Anonymous Authors}

%\iclrfinalcopy
\begin{document}

\maketitle

\begin{abstract}
We study a constructive routing problem for mixture-of-experts layers: given
unlabeled route-relevant features in a spectrally uniform representation space,
can we build a router that dispatches the feature distribution uniformly across
experts without using labels or batch-level capacity forcing?  We define
uniform feature dispatch as an equal-probability partition of feature space and
propose Spectral Partition Routing (SPR), a calibration-based router that
recursively splits unlabeled features by spectral-coordinate quantiles.  SPR can
use only fixed axes, feature-locality objectives, or expert-weight-derived
prototypes; in all cases, quantile thresholds provide the equal-mass guarantee,
while the choice of split direction controls geometry and expert alignment.  We
prove the equal-mass property, clarify why this is not capacity assignment, and
separate the initialization problem from training-time preservation.  The
experiments are left as explicit protocols: static dispatch, non-spherical
spectrally flat features, expert-weight-guided routing, feature-partition
quality, and controlled MoE training dynamics.
\end{abstract}

\section{Introduction}

Sparse mixture-of-experts (MoE) layers route each token representation to a
small subset of experts \citep{shazeer2017outrageously,fedus2022switch}.  A
standard linear top-1 router computes one score per expert and selects the
largest one.  This design is simple, but it leaves the induced feature-space
partition largely determined by random initialization, learned dynamics, and
auxiliary losses.  When the research goal is to understand whether uniformly
represented features can be assigned uniformly to experts, count-level load
balance is not the right starting point.  The primary object is the partition of
feature space itself.

This paper formulates the first step of that problem.  We assume access to an
unlabeled calibration set of route-relevant features in a spectrally uniform
space, meaning that the feature covariance has no dominant direction after the
chosen representation normalization.  We ask whether one can construct a router
that partitions the feature distribution into equal-probability cells, assigns
one cell to each expert, and then routes future tokens by a fixed feature-only
rule.  This setting is intentionally narrower than language-model training: it
does not assume labels, semantic feature identities, or a task-specific notion of
which expert should receive which feature.

The key distinction is between \emph{uniform feature dispatch} and ordinary
load balancing.  Load balancing can be forced inside a finite batch by assigning
tokens to underfilled experts.  Uniform feature dispatch instead asks for a fixed
router $g(x)$ whose induced cells $C_e=\{x:g(x)=e\}$ have equal probability mass
under the feature distribution.  Once calibrated, $g(x)$ must route each token
independently of the other tokens in the current batch.  Under this definition,
using unlabeled feature statistics is allowed; using feature labels, task
targets, or current-batch expert counts is not.

We propose \newterm{Spectral Partition Routing} (SPR).  SPR first standardizes
route-relevant features using unlabeled spectral statistics.  It then constructs
a recursive partition tree.  At each internal node, a split direction is chosen
and a quantile threshold is estimated on the unlabeled features reaching that
node.  The threshold is chosen so that the probability mass sent to the left and
right child matches the number of experts assigned to those children.  Recursing
to $E$ leaves yields $E$ equal-mass feature cells.  The split directions can be
fixed spectral axes, unsupervised locality directions, or directions derived from
expert weights.  Expert weights are important because an expert matrix supplies
an intrinsic candidate prototype through its principal input singular direction,
as in spectral expert-proxy router designs such as MPI \citep{routermpi2026}.
SPR uses such prototypes to choose geometry; quantiles still provide the
uniform-dispatch guarantee.

\paragraph{Contributions.}
\begin{enumerate}
\item We define uniform feature dispatch as an equal-probability feature-space
partition and give a precise non-cheating criterion for routers calibrated from
unlabeled features and optional model-internal expert weights.
\item We introduce Spectral Partition Routing, a practical family of recursive
quantile routers with fixed-axis, locality-driven, and expert-weight-guided
variants.
\item We prove that recursive quantile splitting gives equal-mass expert cells on
the calibration distribution, and we provide an experiment plan that separates
static dispatch, feature locality, expert-prototype alignment, and training-time
preservation.
\end{enumerate}

\section{Problem Definition}

\subsection{Route-relevant feature distribution}

Let $X\in\R^d$ denote the hidden feature to be routed.  In a sequence model,
$X$ should be the representation at the route-relevant position, not an
unfiltered pool of all token positions.  We assume an unlabeled calibration set
\begin{equation}
\calD_{\mathrm{cal}}=\{x_i\}_{i=1}^{n_{\mathrm{cal}}}
\end{equation}
from the same feature population.  No feature identity, sequence label, task
answer, or expert label is available to the method.

The controlled setting is a spectrally uniform representation space.  After
centering by $\mu$ and optionally applying a spectral normalization matrix $B$,
we write
\begin{equation}
  \tilde x = B(x-\mu),
  \qquad
  \mathrm{Cov}(\tilde X)\approx \sigma^2 I.
\end{equation}
If the given features are already spectrally uniform, $B$ may be the identity.
If not, $B$ can be a whitening or shrinkage transform estimated from unlabeled
calibration features.  Spectral uniformity is not treated as proof of semantic
structure.  It only removes the confound that a single high-variance direction
dominates the routing geometry.

\subsection{Uniform feature dispatch}

\begin{definition}[Uniform feature dispatch]
A router is a deterministic map $g:\R^d\to\{1,\ldots,E\}$.  It induces feature
cells
\begin{equation}
 C_e = \{x\in\R^d: g(x)=e\}.
\end{equation}
The router achieves population uniform feature dispatch if
\begin{equation}
 \Prob(X\in C_e)=\Prob(g(X)=e)=\frac{1}{E},\qquad e=1,\ldots,E.
\end{equation}
\end{definition}

This definition is deliberately distributional.  It does not require every
finite batch to contain exactly $N/E$ tokens per expert.  Finite batches should
concentrate around $1/E$ if the calibration distribution matches the deployment
distribution, but exact batch balance is a different mechanism.

\begin{definition}[Non-cheating feature router]
A calibrated router is non-cheating if, after calibration, $g(x)$ depends only on
the current feature $x$, fixed calibration statistics, and optional
model-internal expert weights.  It may not use feature labels, task targets,
evaluation labels, or the current batch's expert counts to adjust assignments.
\end{definition}

Under this definition, using expert weights is allowed.  Expert weights are part
of the model, not data labels.  They can guide how equal-mass feature cells are
associated with experts.  What is not allowed is assigning a token to an expert
because that expert is currently under capacity.

\section{Spectral Partition Routing}

SPR constructs a partition tree.  Each leaf is assigned to one expert.  Each
internal node owns a feature subpopulation and a subset of experts.  A node
splits its feature mass in proportion to how many experts are assigned to each
child.

\subsection{Recursive quantile partition}

Consider a node $\calN$ with feature distribution $P_{\calN}$ and expert subset
$\calE_{\calN}$.  Split the experts into two nonempty groups
$\calE_L$ and $\calE_R$.  Define
\begin{equation}
 \rho_{\calN}=\frac{|\calE_L|}{|\calE_L|+|\calE_R|}.
\end{equation}
Choose a split direction $q_{\calN}\in\Sphere^{d-1}$ and scalar projection
\begin{equation}
 T_{\calN}=q_{\calN}^{\top}\tilde X,\qquad \tilde X\sim P_{\calN}.
\end{equation}
Let $F_{\calN}$ be the CDF of $T_{\calN}$ under the node distribution.  The
threshold is the node quantile
\begin{equation}
 \tau_{\calN}=F_{\calN}^{-1}(\rho_{\calN}).
\end{equation}
The left and right children are
\begin{align}
 C_L &= \{x\in C_{\calN}: q_{\calN}^{\top}\tilde x\le \tau_{\calN}\},\\
 C_R &= \{x\in C_{\calN}: q_{\calN}^{\top}\tilde x> \tau_{\calN}\}.
\end{align}
The recursion stops when a node contains one expert.  That leaf cell is assigned
to that expert.  If $E=2^K$, all splits can be medians.  If $E$ is not a power
of two, the split ratio follows the left/right expert counts.

\subsection{Choosing the split direction}

The quantile threshold gives equal mass.  The split direction controls geometry,
locality, and expert alignment.  We study three variants.

\paragraph{SPR-Fixed.}
A fixed sequence of directions is chosen before seeing labels or tasks.  Examples
include a fixed random orthogonal basis, a Hadamard basis, or a Fourier-like
basis.  This is the cleanest construction: it uses only unlabeled feature
quantiles and has no expert dependence.

\paragraph{SPR-Local.}
The direction is chosen from unlabeled calibration features to improve geometric
locality.  A practical implementation samples candidate unit directions
$\{q^{(m)}\}_{m=1}^M$, computes the quantile split for each, and chooses the one
with lowest within-cell variance,
\begin{equation}
  \mathcal{L}_{\mathrm{local}}(q)
  =\sum_{b\in\{L,R\}} \Prob(C_b\mid C_{\calN})
  \E\left[\|\tilde X-\mu_b\|_2^2\mid \tilde X\in C_b\right].
\end{equation}
This remains unsupervised.  It uses feature geometry, not labels.

\paragraph{SPR-Expert.}
Expert weights can guide the partition.  Let $W_e$ be an input-side expert
matrix.  Define an expert spectral prototype
\begin{equation}
 v_e\in\arg\max_{\|v\|_2=1}\|v^{\top}W_e\|_2^2,
\end{equation}
which is a top eigenvector of $W_eW_e^{\top}$.  This matches the expert-proxy
view of spectral router designs: MPI-style updates move router rows toward
principal expert directions through power iteration and retraction
\citep{routermpi2026}.

At node $\calN$, cluster or split the expert prototypes
$\{v_e:e\in\calE_{\calN}\}$ into $\calE_L$ and $\calE_R$.  Let
\begin{equation}
 \bar v_L=\frac{1}{|\calE_L|}\sum_{e\in\calE_L} v_e,\qquad
 \bar v_R=\frac{1}{|\calE_R|}\sum_{e\in\calE_R} v_e.
\end{equation}
Then use
\begin{equation}
 q_{\calN}=\frac{\bar v_L-\bar v_R}{\|\bar v_L-\bar v_R\|_2}
\end{equation}
whenever the numerator is nonzero; otherwise fall back to SPR-Fixed or
SPR-Local.  Expert weights therefore influence which equal-mass cell goes to
which expert.  They do not replace the quantile rule that enforces uniform
feature dispatch.

\subsection{Router forms for training}

SPR can be used in three training modes.

\paragraph{Frozen-SPR.}
The tree, thresholds, and expert assignment are fixed after calibration.  Only
experts and the rest of the model train.  This isolates whether a uniform feature
partition is useful when it is preserved exactly.

\paragraph{SPR-init-linear.}
SPR is used only to initialize a conventional linear router.  Let
\begin{equation}
 c_e=\E[\tilde X\mid g_{\mathrm{SPR}}(X)=e]
\end{equation}
be the cell center.  The linear router row for expert $e$ is initialized to
$c_e/\|c_e\|_2$ when the center is nonzero.  This tests whether the learned
linear top-1 router can preserve a feature partition discovered by SPR.

\paragraph{SlowUpdate-SPR.}
The tree is recalibrated periodically using a moving window of unlabeled
features.  Updates are infrequent and smoothed to avoid boundary oscillation.
This mode studies representation drift while keeping the router non-cheating:
each route decision still uses only the current feature and fixed statistics from
previous calibration windows.

\section{Theory}

\subsection{A quantile split gives prescribed mass}

\begin{theorem}[Node quantile split]
\label{thm:quantile_split}
Let $X\sim P_{\calN}$ be the feature distribution at node $\calN$, let
$q\in\Sphere^{d-1}$, and let $T=q^{\top}\tilde X$ have a continuous CDF $F$.
For $\rho\in(0,1)$, define $\tau=F^{-1}(\rho)$ and split by
$T\le \tau$.  Then
\begin{equation}
 \Prob(T\le \tau)=\rho,\qquad \Prob(T>\tau)=1-\rho.
\end{equation}
\end{theorem}

\begin{proof}
This is the defining property of the $\rho$-quantile for a continuous scalar
random variable.
\end{proof}

\subsection{Recursive splitting gives uniform expert cells}

\begin{theorem}[Uniform dispatch by recursive spectral partition]
\label{thm:recursive_uniform}
Suppose every internal node $\calN$ in the SPR tree splits its feature mass in
proportion to the number of experts assigned to its children:
\begin{equation}
 \Prob(C_L\mid C_{\calN})=\frac{|\calE_L|}{|\calE_{\calN}|},\qquad
 \Prob(C_R\mid C_{\calN})=\frac{|\calE_R|}{|\calE_{\calN}|}.
\end{equation}
If the root contains all $E$ experts and each leaf contains one expert, then each
leaf cell $C_e$ has mass $1/E$.
\end{theorem}

\begin{proof}
For any node $\calN$, we prove by induction that
\begin{equation}
 \Prob(C_{\calN})=\frac{|\calE_{\calN}|}{E}.
\end{equation}
The root satisfies this because its mass is $1$ and it contains $E$ experts.  If
the claim holds for a parent node, then for its left child,
\begin{equation}
 \Prob(C_L)=\Prob(C_{\calN})\Prob(C_L\mid C_{\calN})
 =\frac{|\calE_{\calN}|}{E}\frac{|\calE_L|}{|\calE_{\calN}|}
 =\frac{|\calE_L|}{E}.
\end{equation}
The right child is identical.  At a leaf, $|\calE_e|=1$, so
$\Prob(C_e)=1/E$.
\end{proof}

\subsection{Why this is not capacity assignment}

\begin{proposition}[Fixed feature router]
After calibration, SPR defines a deterministic function $g(x)$ of the current
feature and stored tree parameters.  It does not depend on the other examples in
the current batch.  Therefore SPR is not a batch-level capacity assignment.
\end{proposition}

\begin{proof}
The route of $x$ is computed by evaluating the stored split directions and
thresholds along one root-to-leaf path.  No term in this computation contains the
current batch composition or the current number of examples already assigned to
any expert.
\end{proof}

\subsection{Finite calibration and deployment error}

In practice, $F_{\calN}$ is estimated from a finite calibration set.  Empirical
quantiles give exact or near-exact balance on the calibration set, but holdout
usage depends on calibration size and distribution shift.  For one scalar split,
standard empirical CDF concentration implies a deviation scale
$O(\sqrt{\log(1/\delta)/n_{\calN}})$ at node sample size $n_{\calN}$
\citep{dvoretzky1956asymptotic}.  In a tree, lower nodes have fewer samples, so
calibration size is a real experimental variable.  This is a statistical
estimation issue, not a capacity-forcing mechanism.

\subsection{Role of expert weights}

Expert weights affect the semantics of the partition, not the equal-mass proof.
If SPR-Expert chooses $q_{\calN}$ from expert prototypes, Theorem
\ref{thm:quantile_split} still guarantees the requested node mass because the
threshold is a quantile of the feature projection.  Poor expert prototypes can
produce geometrically poor or low-margin cells, but they cannot by themselves
break the population mass guarantee when the quantiles are calibrated correctly.
This separation is useful: expert weights provide an alignment prior, while
quantiles provide dispatch uniformity.

\section{Training Dynamics}

Uniform dispatch at calibration time is only the initialization problem.  During
training, the encoder, router statistics, and expert weights may move.  A top-1
router is especially sensitive because a small boundary crossing changes the
selected expert discontinuously.

For tree routers, define the route margin of $x$ as the closest normalized
distance to a threshold along the selected path:
\begin{equation}
 m_{\mathrm{tree}}(x)=\min_{\calN\in\mathrm{path}(x)}
 \frac{|q_{\calN}^{\top}\tilde x-\tau_{\calN}|}{\|q_{\calN}\|_2}.
\end{equation}
For linear top-1 routers, use the score gap
\begin{equation}
 m_{\mathrm{linear}}(x)=s_{a(x)}(x)-\max_{j\ne a(x)}s_j(x).
\end{equation}
Small lower-tail margins predict route flips under representation drift.  The
training experiments should therefore report not only loss and expert usage, but
also route flip rate, feature-to-route agreement, and margin lower tails.

The central training hypothesis is modest: a calibrated feature partition can
provide a better initialization for controlled uniform-feature tasks than a
random linear router, but preservation requires sufficient margin or slow
recalibration.  This paper does not assume SPR automatically discovers semantic
experts; it tests whether a clean equal-mass feature partition can be constructed
and whether it survives controlled optimization.

\section{Experiments}
\label{sec:experiments}

This section is intentionally specified as a protocol.  No experimental results
are claimed in this draft.  Each experiment is designed to test one part of the
story.

\subsection{Static uniform dispatch}

\paragraph{Goal.}
Verify that SPR constructs equal-mass feature cells on spectrally uniform data
and generalizes from calibration to holdout samples.

\paragraph{Data.}
Use Gaussian features $X\sim\mathcal{N}(0,I_d)$ and normalized spherical
features $X/\|X\|_2\sim\mathrm{Unif}(\Sphere^{d-1})$.  Sweep $d$, $E$, and
calibration size.

\paragraph{Methods.}
Compare vanilla random linear top-1 routing, regular-simplex prototype routing,
SPR-Fixed, SPR-Local, SPR-Expert with synthetic expert weights, and balanced
assignment as a forced-count control.

\paragraph{Metrics.}
Report calibration LoadCV, holdout LoadCV, maximum absolute usage error, margin
lower tail, and calibration-to-holdout gap.

\subsection{Spectrally flat but non-spherical features}

\paragraph{Goal.}
Show that covariance flatness alone does not make fixed prototype routers
reliable, while quantile-calibrated SPR still constructs equal-mass cells.

\paragraph{Data.}
Use axis-supported distributions, mixtures of symmetric directions, and other
zero-mean distributions with covariance close to $I/d$ but non-spherical support.

\paragraph{Methods and expectation.}
Prototype-only routers may be imbalanced because their cells receive unequal
probability mass.  SPR should remain balanced on the calibration distribution and
near-balanced on holdout data if the distribution is stable.

\subsection{Expert-weight-guided partitioning}

\paragraph{Goal.}
Audit the role of expert weights.  Expert spectral directions should help align
feature cells to experts, but quantile thresholds should be responsible for mass
uniformity.

\paragraph{Data and experts.}
Use both standard randomly initialized expert matrices and controlled expert
matrices with prescribed principal directions.  Extract expert prototypes by SVD
or power iteration.

\paragraph{Methods.}
Compare expert-SVD prototype top-1 routing, MPI-style one-step prototypes,
SPR-Expert, SPR-Fixed, and SPR-Local.  In SPR-Expert, expert prototypes choose
split directions and expert groupings; quantiles choose thresholds.

\paragraph{Metrics.}
Report expert-prototype spread, prototype alignment, LoadCV, margin, and the
agreement between SPR leaves and expert-prototype groups.

\subsection{Feature partition quality}

\paragraph{Goal.}
Distinguish equal-mass dispatch from meaningful feature partition.

\paragraph{Data.}
Construct synthetic feature classes $F$ for evaluation only.  The router never
uses $F$.  Generate features from class-conditioned distributions, then whiten or
otherwise normalize them so that global covariance is close to identity.

\paragraph{Metrics.}
Report LoadCV, route--feature normalized mutual information, expert purity,
within-cell variance, and margin lower tail.  This experiment tests whether SPR
cuts geometry in a way that preserves latent feature groups, not merely counts.

\subsection{Controlled MoE training dynamics}

\paragraph{Goal.}
Test whether a calibrated feature partition helps controlled MoE training and
whether it is preserved under optimization.

\paragraph{Task.}
Construct a toy supervised task where each latent feature class requires a
different function, for example $Y=A_FX+\epsilon$.  The router methods cannot use
$F$ during calibration or training, but evaluation can measure whether routes
align with $F$.

\paragraph{Routers.}
Compare vanilla linear top-1, linear top-1 with load-balancing loss,
regular-simplex initialization, Frozen-SPR, SPR-init-linear, SlowUpdate-SPR, and
SPR-Expert.

\paragraph{Metrics.}
Report train and validation loss, expert usage, route--feature NMI, route flip
rate, margin p05, expert specialization score, and representation drift.

\subsection{Ablations}

Ablate calibration size, feature dimension, expert count, non-power-of-two
expert counts, split-direction policy, whitening on/off, expert-weight quality,
slow-update frequency, and feature drift.  Balanced assignment should remain only
as a control for count forcing, not as the target method.

\section{Discussion}

The paper's target is intentionally first-stage.  Uniform feature dispatch means
that the feature distribution is partitioned into equal-probability expert cells.
It does not mean that the cells are semantically optimal, nor that the resulting
experts improve downstream performance.  Those are later questions that require
feature-structure and task-loss experiments.

The main design tension is the choice of split direction.  Quantiles guarantee
mass once a direction is chosen.  Fixed axes are clean but may ignore local
feature geometry.  Locality-driven directions are more adaptive but require more
calibration.  Expert-weight-guided directions connect the feature partition to
the MoE parameters.  This last route is important: if an expert matrix already
contains a strong principal direction, using it to guide a feature split is a
reasonable way to tie router cells to expert identity.  However, expert weights
alone do not guarantee uniform dispatch; they must be paired with calibrated
thresholds.

Training dynamics remain a separate issue.  A good initial partition can still
be lost if representations drift across thresholds or if a learnable top-1
router collapses into early lock-in.  This is why the second-stage experiments
must report margin and route flips, not only usage counts.  In controlled
sequence or slot-feature settings, the first practical question is whether the
route-relevant states are themselves spectrally well behaved.  If all positions
are pooled together, route and non-route states may mix, making the calibration
population ill-defined.  The correct workflow is therefore: identify the
route-relevant feature population, construct an equal-mass feature partition,
attach the cells to experts, and then test whether training preserves the
partition.

\section{Scope and Limitations}

SPR is a construction for unlabeled feature dispatch, not a proof that any
standard MoE router will discover the same partition.  It assumes a calibration
feature population representative of later routed features.  If the feature
distribution shifts, stored quantiles can become stale.  If the distribution has
large atoms, no deterministic feature-only router can split those identical
features across experts without either stochastic tie-breaking or violating the
principle that identical features receive identical routes.  Finally, equal-mass
cells are not equivalent to useful specialization.  Feature quality and training
utility must be measured separately.

\section{Conclusion}

We reformulate uniform routing as a constructive feature-partition problem.  In a
spectrally uniform feature setting, the goal is to build a fixed, non-cheating
router whose cells have equal probability mass.  Spectral Partition Routing
achieves this by combining unlabeled spectral calibration, recursive quantile
splitting, and optional expert-weight-guided directions.  Quantiles provide the
uniform-dispatch guarantee; split-direction policy provides geometry and expert
alignment.  The resulting story separates three questions that are often mixed:
can we construct an equal-mass feature partition, can that partition align with
expert identity, and can it survive training dynamics?  This draft resolves the
first question by construction and lays out the experiments needed for the next
two.

\bibliographystyle{iclr2026_conference}
\bibliography{references}

\end{document}
